\section{Introduction}
\label{sec:intro}

The value of additional cases in reproduction-number surveillance depends on the quantity being estimated. A current transmission level, the average of transmission levels during an observed interval, and the mean of a process generating successive intervals are different inferential targets. Each can be relevant to surveillance, but they assign different roles to variation between periods. A change in transmission is part of the signal when estimating the current value; it contributes sampling variation when estimating a process mean. Statements about a precision limit therefore require both a target and an observation model. Case volume alone cannot determine that limit.

Reproduction numbers support epidemic assessment by summarising transmission under specified conditions \citep{fraser2007estimating, gostic2020practical}. Renewal methods relate incidence to past infectiousness through a generation-interval distribution, and their uncertainty depends on the information in the observed counts \citep{cori2013new, thompson2019improved}. Work on low incidence, reporting noise and timely detection has clarified important limits of period-specific inference \citep{parag2021improved, parag2022quantifying, parag2025asymmetric}. Temporal aggregation and delays create further differences between transmission and the apparent reproduction number of an observation series \citep{nash2023estimating, ogigittins2024simulation, eales2024differences}. These results motivate a distinction between uncertainty about a realised transmission trajectory and uncertainty about a population of periods from which that trajectory is sampled.

We develop that distinction in a one-generation branching experiment. Primary cases have independent negative binomial offspring, each secondary infection is reported independently, and the number of primary cases is fixed or exogenous. With a fixed, known reporting probability, a scaled count is conditionally unbiased for the current reproduction number. Its expected squared error decreases inversely with case volume. Averaging over a fixed set of observed periods similarly reduces counting error when the target is their realised mean. Estimating a stationary process mean adds the sampling variance of those periods. Increasing cases within the same periods leaves this additional term unchanged, whereas observing further periods can reduce it according to their dependence structure.

This distinction builds on a classical variance structure. Population ecology separates independent individual outcomes from environmental variation shared across a population \citep{engen1998demographic, lande2003stochastic}; stochastic epidemic models likewise distinguish demographic and environmental contributions \citep{breto2009time, he2010plug}. Shared conditions can include climate, contact patterns and school calendars \citep{cauchemez2008estimating, shaman2009absolute}. In sampling theory, additional observations within a cluster reduce within-cluster error, while additional clusters address between-cluster variation \citep{cochran1977sampling, hayes1999simple}. The contribution here is to specify the analogous terms for reproduction-number surveillance and show how changing the target changes the relevant risk. The resulting balance scale compares two contributions under a stated model; it is not a universal minimum case requirement.

Reporting requires separate attention. Independent omission of infections differs from a reporting probability that changes for all infections in the same period \citep{azmon2014estimating, stoner2019hierarchical}. Even a stable reporting probability can produce bias if its value is unknown or the count is incorrectly normalised. When reporting probabilities fluctuate independently of transmission, normalisation by their known mean leaves an additional error term for current-value estimation. It also affects the realised finite-period mean. For process-mean estimation, count means and variances alone generally cannot separate reporting volatility from transmission variation. Such ambiguity is consistent with the broader difficulty of separating process and observation variation in time series \citep{knape2008estimability, augermethe2016state}.

Routine hospital surveillance does not directly observe the variables of the branching experiment. Admissions follow infection after a delay, aggregate heterogeneous populations and are recorded on a calendar schedule. Their count scales are outcomes of the epidemic, rather than exogenous design variables. We therefore examine squared relative deviations of apparent weekly growth statistics in US National Healthcare Safety Network influenza admissions \citep{cdc2026nhsn, frutos2025evaluation}. These deviations describe variation around an empirical centre; they are not squared errors against observed transmission truth. An affine association with inverse count can describe this variation without identifying its biological sources or establishing an asymptote outside the observed count range.

The empirical analyses distinguish questions that the data can address. Comparing affine and power-law curves concerns descriptive shape and prediction. Controlled binomial thinning examines the response to additional sampling of an observed trajectory. State-intercept fits examine within-state associations, while alternative centres assess sensitivity to common temporal patterns. Simulation connects these summaries to specified transmission and observation mechanisms, conditional on generation intervals, delays and dependence. None of these operations alone identifies environmental transmission variance. Their roles, uncertainty calculations and execution status are documented in Methods and Results (Sections~\ref{sec:methods} and~\ref{sec:results}).

The paper thus places inferential targets before case-volume interpretations. Section~\ref{sec:background} states the exact identities and the moment nonidentification result, with the targets collected in Table~\ref{tab:targets}. The empirical analysis then assesses associations and their sensitivity on the scale of the observed statistics. Proofs and symbolic extensions for forecasting, shrinkage and design appear in the appendix. Together, these components clarify when additional cases improve the chosen estimate and which remaining uncertainties require information beyond case counts.
